\textbf{COLUNI 2021 -- Questão 10}\par\medskip Analise a figura de linguagem – personificação – em destaque no excerto:

\par

“\underline{Meu coração pinoteou de alvoroço} e eu disparei de volta à cabina, tropeçando nos marinheiros [...]”. (LAURITO, Ilka Brunhilde. \textbf{A Menina que fez a América}. São Paulo, FTD, 2019. p. 86)

\par

Considere as afirmativas que apresentam a mesma figura de linguagem identificada no excerto, atribuindo V para a(s) verdadeira(s) e F para a(s) falsa(s).

\par

( ) “Suas mãos, grandes, quentes, calosas, mãos de quem mexia na terra, \underline{mãos enraizadas no solo}.” (LAURITO, Ilka Brunhilde. \textbf{A Menina que fez a América}. São Paulo, FTD, 2019. p. 11) ( ) “Aí ele falava que \underline{as estrelas}, como as pessoas aqui da terra, \underline{também não gostavam de viver sozinhas}.” (LAURITO, Ilka Brunhilde. \textbf{A Menina que fez a América}. São Paulo, FTD, 2019. p. 14) ( ) “E a noite ficou presa entre os meus dedos, \underline{silenciosamente adormecida} até o alvorecer.” (LAURITO, Ilka Brunhilde. \textbf{A Menina que fez a América}. São Paulo, FTD, 2019. p. 31) ( ) “Escancarei-a e a janelinha \underline{se debruçou} sobre os telhados da casa de Rina.” (LAURITO, Ilka Brunhilde. \textbf{A Menina que fez a América}. São Paulo, FTD, 2019. p. 39)

\par

Assinale a sequência CORRETA:\begin{enumerate}
\item V, F, F, V;
\item V, F, F, F;
\item F, V, V, V;
\item F, V, V, F.

\end{enumerate}
